\section{\Zoo{} 3.0 \GPU{}-Native Substrate}
\label{sec:zoo-3-0-substrate}

The per-\LLM{} chain framework specified in this paper executes on the
\Zoo{} 3.0 \GPU{}-native substrate. This section anchors the
relationship between the per-\LLM{} chain abstraction and the underlying
\GPU{} primitives that make it economically viable.

\subsection{\GPU{}-residency under per-\LLM{} chain workloads}

A per-\LLM{} chain workload is a constant stream of attestation
verifications, attribution-graph updates, and inference-receipt
rollups. Each of these operations is dominated by signature
verification, hash construction, and small-state mutation. Running
this stream on a CPU \textsc{evm} bottlenecks at the
verification step. \Zoo{} 3.0 takes the entire stream onto the
\GPU{} per the \GPU{}-Residency Invariant
(\textsc{lp-137})~\cite{lp-137-residency}: data lands on the device
when an attestation enters the mempool and stays on the device through
verification, ledger mutation, and certificate emission.

\subsection{\QuasarGPU{} adapter (\textsc{lp-132})}

The \QuasarGPU{} adapter~\cite{lp-132-quasargpu} schedules deterministic
\GPU{} kernels for triple-consensus signature verification. For per-\LLM{}
chains this matters because the chain emits a high-volume of small,
heterogeneous transactions (training receipt, eval receipt, attribution
update) that the adapter batches into a single device-side verification
pass per block. The result is that a per-\LLM{} chain block clears in
the same wall-clock time as a generic \Zoo{} block despite the different
transaction mix.

\subsection{Cross-reference}

The full \Zoo{} 3.0 substrate specification, including the L1 graduation
mechanics, native \DEX{}, AI desktop, bot, bridge, and privacy paths,
is in the \emph{Zoo 3.0 Launch} paper~\cite{zoo-3-0-launch}. Per-\LLM{}
chains are the canonical first-class workload of that substrate; the
two specifications activate together on 2025-12-25.
